\documentclass{article}
\usepackage[utf8]{inputenc}
\usepackage{amsmath}
\usepackage{amssymb}
\usepackage{graphicx}
\usepackage{booktabs}
\usepackage{float}
\usepackage{newtxtext}
\usepackage{hyperref}
\usepackage{algorithm}
\usepackage{algpseudocode}

\title{Comparative Analysis of Machine Learning Approaches for MNIST
Handwritten Digit Recognition}
\author{CheesyChocolate}
\date{\today}

\begin{document}

\maketitle

\begin{abstract}
  This paper presents a comprehensive comparison of multiple machine
  learning approaches for handwritten digit recognition using the
  MNIST dataset. We implement and evaluate three different models:
  Support Vector Machine (SVM), Random Forest, and K-Nearest
  Neighbors (KNN). Each model is evaluated using standardized
  metrics, and their strengths and weaknesses are analyzed in the
  context of pattern recognition. Our findings provide insights into
  the effectiveness of different algorithmic approaches for this
  classic computer vision task.
\end{abstract}

\section{Introduction}
Handwritten digit recognition represents a fundamental pattern
recognition problem in computer vision and machine learning. The
MNIST dataset, consisting of 60,000 training images and 10,000 test
images of handwritten digits (0-9), has become a standard benchmark
for evaluating classification algorithms. Despite its apparent
simplicity, the task involves significant challenges including
variations in writing styles, stroke thickness, and digit positioning.

This study compares three distinct machine learning approaches to
solve the MNIST classification problem. By implementing these models
using a consistent methodology and evaluation framework, we aim to
provide a clear understanding of their relative strengths,
weaknesses, and appropriate use cases.

\begin{figure}[ht]
  \centering
  \includegraphics[width=0.8\textwidth]{fig/mnist_samples.png}
  \caption{Sample digits from the MNIST dataset (0-9)}
  \label{fig:mnist_samples}
\end{figure}

\section{Related Work}
The MNIST dataset has been extensively studied since its introduction
by LeCun et al.. Early approaches to the problem utilized techniques
such as linear classifiers, k-nearest neighbors, and support vector
machines, achieving accuracies in the range of 84\% to 98\%.

With the advent of deep learning, convolutional neural networks
(CNNs) have demonstrated superior performance on this task. LeNet-5,
one of the earliest CNN architectures, achieved an error rate of
approximately 0.8\%. More recent architectures have further reduced
the error rate, with some sophisticated models reaching error rates below 0.2\%.

Despite the superior performance of deep learning approaches,
traditional machine learning algorithms remain valuable for their
interpretability, reduced computational requirements, and
effectiveness in scenarios with limited training data.

\section{Methodology}

\subsection{Dataset}
The MNIST dataset consists of 28$\times$28 pixel grayscale images of
handwritten digits. The dataset is divided into 60,000 training
images and 10,000 test images. Each pixel in the images has an
intensity value ranging from 0 to 255, which we normalize to the
range [0,1]. For our experiments, we used a randomly selected subset
of 5,000 images to reduce computational requirements, with 80\% for
training and 20\% for testing.

\subsection{Models}
We implement and evaluate three different models:

\subsubsection{Support Vector Machine (SVM)}
We utilize an SVM with an RBF kernel, which can capture non-linear
relationships in the data. The hyperparameters C (regularization
parameter) and gamma (kernel coefficient) are set to 1.0 and 'scale'
respectively.

\subsubsection{Random Forest}
Our Random Forest classifier consists of 100 decision trees. The
model leverages ensemble learning to improve accuracy and reduce overfitting.

\subsubsection{K-Nearest Neighbors (KNN)}
The KNN classifier is implemented with k=5 neighbors and uniform
weighting. Despite its simplicity, KNN can be effective for the MNIST
task due to the clustering of similar digits in the feature space.

\subsection{Evaluation Metrics}
We evaluate each model using the following metrics:
\begin{itemize}
  \item Accuracy: The proportion of correctly classified digits
  \item Precision: The proportion of true positives among all
    positive predictions
  \item Recall: The proportion of true positives among all actual positives
  \item F1 Score: The harmonic mean of precision and recall
  \item Confusion Matrix: A visual representation of the model's
    classification performance
\end{itemize}

\section{Results}
\subsection{Model Performance}
The performance comparison of the three models is shown in Figure
\ref{fig:model_comparison}. The SVM model achieved the highest
accuracy at 97.2\%, followed by Random Forest at 96.9\% and KNN at
96.7\%. While the differences in accuracy are relatively small, they
can be significant when applied to large-scale real-world applications.

\begin{figure}[ht]
  \centering
  \includegraphics[width=0.7\textwidth]{fig/model_comparison.png}
  \caption{Accuracy comparison of the three models}
  \label{fig:model_comparison}
\end{figure}

Each model demonstrates different strengths in classifying specific
digits. The precision, recall, and F1-score for each model across the
ten digit classes are summarized in Tables \ref{tab:svm_metrics},
\ref{tab:rf_metrics}, and \ref{tab:knn_metrics}.

\subsection{Confusion Matrices}
The confusion matrices provide insights into the specific patterns of
misclassifications for each model. Figure
\ref{fig:confusion_matrices} shows the confusion matrices for all three models.

\begin{figure}[ht]
  \centering
  \includegraphics[width=0.49\textwidth]{fig/SVM_confusion_matrix.png}
  \includegraphics[width=0.49\textwidth]{fig/RandomForest_confusion_matrix.png}
  \includegraphics[width=0.49\textwidth]{fig/KNN_confusion_matrix.png}
  \caption{Confusion matrices for SVM (top left), Random Forest (top
  right), and KNN (bottom)}
  \label{fig:confusion_matrices}
\end{figure}

Common patterns of misclassification across all models include:
\begin{itemize}
  \item Confusion between 4 and 9, due to similar upper loops
  \item Confusion between 3 and 5, which differ mainly in the
    orientation of strokes
  \item Confusion between 7 and 1, especially when the 7 is written
    without the horizontal stroke
\end{itemize}

\subsection{Sample Predictions}
Figure \ref{fig:sample_predictions} shows sample predictions from the
SVM model on test data. These examples illustrate both correct
classifications and typical errors.

\begin{figure}[ht]
  \centering
  \includegraphics[width=0.95\textwidth]{fig/SVM_sample_predictions.png}
  \caption{Sample predictions from the SVM model}
  \label{fig:sample_predictions}
\end{figure}

\section{Discussion}
\subsection{Comparison of Approaches}
Each of the three models has distinct characteristics that make them
suitable for different scenarios:

\subsubsection{Support Vector Machine}
The SVM model achieved the highest overall accuracy. SVMs excel at
finding optimal decision boundaries in high-dimensional spaces,
making them well-suited for image classification tasks. The RBF
kernel effectively captures the non-linear relationships in the data.
However, SVMs have higher computational complexity during training,
especially for large datasets.

\subsubsection{Random Forest}
The Random Forest model performed nearly as well as the SVM, with the
advantage of being more interpretable and faster to train. The
ensemble nature of Random Forests makes them robust against
overfitting, and they can naturally handle multi-class problems.
Random Forests also provide insights into feature importance, which
can be valuable for understanding the classification process.

\subsubsection{K-Nearest Neighbors}
The KNN model, despite its simplicity, achieved competitive results.
KNN is easy to implement and understand, with no training phase
required. However, it has higher memory requirements as it needs to
store the entire training set, and prediction can be slower as the
dataset grows. KNN is particularly sensitive to the curse of
dimensionality, which may limit its effectiveness on very high-dimensional data.

\subsection{Practical Considerations}
Beyond accuracy, several practical considerations influence the
choice of model for real-world applications:

\begin{itemize}
  \item \textbf{Training time:} Random Forest and SVM require a
    significant training phase, while KNN has no training phase but
    higher inference costs.
  \item \textbf{Memory usage:} KNN requires storing the entire
    training set, making it memory-intensive for large datasets.
  \item \textbf{Interpretability:} Random Forest offers better
    interpretability through feature importance measures compared to
    SVM and KNN.
  \item \textbf{Scalability:} SVM and KNN do not scale as well to
    very large datasets compared to Random Forest, which can be
    trained in parallel.
\end{itemize}

\section{Conclusion}
This study has demonstrated that traditional machine learning
approaches can achieve high accuracy on the MNIST handwritten digit
recognition task. The SVM model achieved the best performance,
closely followed by Random Forest and KNN. Each model offers
different trade-offs in terms of accuracy, training time, memory
requirements, and interpretability.

While deep learning approaches like CNNs can achieve even higher
accuracy on this task, traditional machine learning methods remain
relevant due to their lower computational requirements, faster
training, and better interpretability. They serve as strong baselines
and are often sufficient for many practical applications.

The choice of model should be guided by the specific requirements of
the application, including the available computational resources, the
need for interpretability, and the size of the dataset.

\section{Future Work}
Future directions for this research include:

\begin{itemize}
  \item Hyperparameter optimization to further improve model performance
  \item Exploration of feature engineering techniques to enhance the
    representation of the digits
  \item Comparison with other machine learning approaches such as
    decision trees and logistic regression
  \item Investigation of ensemble methods combining multiple model types
  \item Application to more complex datasets beyond MNIST
\end{itemize}

\end{document}
